\chapter{Endocrine Disorders}

\section*{Patient focus}
Endocrine glands produce hormones that control many body functions. Imaging helps find tumours, cysts, or enlargement of these glands. For the thyroid, an ultrasound is the main test and can determine whether a lump needs a biopsy. For the adrenal glands, a CT scan can tell whether a mass is a benign adenoma or something more serious. For the pituitary gland in the brain, an MRI is the best test. Your doctor will choose the right imaging based on which hormone problem is suspected.


\label{ch:ch66_endocrine_disorders}

\section{Imaging}
Endocrine imaging targets the thyroid, parathyroids, adrenal glands, pituitary, and neuroendocrine tumors. Ultrasound, CT, MRI, nuclear medicine, and PET are selected by anatomy, hormone findings, and the clinical question.

\section{Thyroid Imaging}

\subsection{Thyroid Ultrasound and TI-RADS}
Ultrasound is the first-line investigation for all thyroid nodules. The \textbf{ACR TI-RADS} system standardises risk stratification:

\begin{table}[h]
\centering
\caption{ACR TI-RADS scoring summary}
\begin{tabular}{|l|l|}
\hline
\textbf{Feature} & \textbf{Points} \\
\hline
\textbf{Composition} & \\
Cystic or spongiform & 0 pts \\
\textbf{Echogenic foci} & \\
\textbf{Echogenicity} & \\
\hline
Hypoechoic & 2 pts \\
Isoechoic or hyperechoic & 1 pt \\
Lobulated & 2 pts \\
Macrocalcifications & 2 pts \\
\hline
\textbf{Margin} & \\
Microcalcifications & 3 pts \\
Mostly solid & 1 pt \\
\hline
\textbf{Shape} & \\
Smooth or ill-defined & 0 pts \\
Solid & 2 pts \\
Spiculated or irregular & 3 pts \\
\hline
Taller than wide & 3 pts \\
Very hypoechoic & 3 pts \\
Wider than tall & 0 pts \\
\hline
\end{tabular}
\end{table}

\begin{itemize}
  \item \textbf{TI-RADS 1} (0 points): Benign---no FNA
  \item \textbf{TI-RADS 2} (2 points): Not suspicious---no FNA
  \item \textbf{TI-RADS 3} (3 points): Mildly suspicious---FNA if $\geq$2.5 cm
  \item \textbf{TI-RADS 4} (4--6 points): Moderately suspicious---FNA if $\geq$1.5 cm
  \item \textbf{TI-RADS 5} ($\geq$7 points): Highly suspicious---FNA if $\geq$1 cm
\end{itemize}

\begin{pearlbox}[TI-RADS and FNA]
TI-RADS determines whether fine needle aspiration is recommended and at what size threshold. A suspicious 1 cm nodule may need FNA, while a benign-appearing 2 cm nodule can be safely followed. Always correlate with clinical risk factors (radiation history, family history, personal history of head/neck radiation).
\end{pearlbox}

\section{Parathyroid Imaging}
\begin{itemize}
  \item \textbf{Indication}: localising parathyroid adenoma before surgery in primary hyperparathyroidism
  \item \textbf{US}: hypoechogenic, oval mass posterior to thyroid; sensitivity 70--80\%
  \item \textbf{Sestamibi (Tc-99m)}: retained tracer in adenoma at 2--3 hours (thyroid washes out); sensitivity 80--90\%
  \item \textbf{4D CT}: three contrast-enhanced phases to localise adenoma; sensitivity $>$90\%; useful when US and sestamibi are discordant
  \item \textbf{SPECT/CT}: fused nuclear and CT imaging for precise localisation, especially in ectopic or mediastinal adenomas
\end{itemize}

\section{Adrenal Imaging}

\subsection{Adrenal Incidentaloma Workup}
Adrenal incidentalomas are found in $\sim$5\% of abdominal CT scans. Characterisation is essential:

\begin{table}[h]
\centering
\caption{Adrenal adenoma washout CT criteria}
\begin{tabular}{|l|l|}
\hline
\textbf{Parameter} & \textbf{Value} \\
\hline
Absolute washout (30 min) & $>$60\% \\
Relative washout (15 min) & $>$40\% \\
Size & Typically $<$4 cm, $<$10 g \\
Unenhanced attenuation & $\leq$10 HU (lipid-rich adenoma) \\
\hline
\end{tabular}
\end{table}

\begin{itemize}
  \item \textbf{Lipid-rich adenoma} ($\sim$70\%): unenhanced attenuation $\leq$10 HU---diagnostic without further imaging
  \item \textbf{Lipid-poor adenoma}: higher unenhanced attenuation; requires washout CT or MRI chemical shift
  \item \textbf{MRI chemical shift}: signal drop on opposed-phase images indicates intracellular lipid---diagnostic of adenoma
  \item \textbf{Pheochromocytoma}: avid enhancement, may be cystic or necrotic, unenhanced attenuation usually $>$10 HU, washout $<$60\%
  \item \textbf{Adrenal carcinoma}: large ($>$4 cm), irregular, heterogeneous, necrotic, may have calcification
  \item \textbf{Metastasis}: solid, enhancing, may be bilateral---known primary cancer is key context
\end{itemize}

\begin{pearlbox}[Pheochromocytoma always rule out]
Before any biopsy or surgical intervention on an adrenal mass, biochemical exclusion of phaeochromocytoma is mandatory. Unrecognised phaeochromocytoma can cause hypertensive crisis during surgery. Always check plasma metanephrines.
\end{pearlbox}

\section{Pituitary Imaging}

\subsection{Pituitary Protocol MRI}
\begin{itemize}
  \item \textbf{Sequences}: thin-coronal T1 (pre and post gadolinium), sagittal T1, coronal T2, dynamic post-contrast T1, axial DWI
  \item \textbf{Microadenoma} ($<$10 mm): hypointense on T1, may show delayed enhancement compared to normal pituitary
  \item \textbf{Macroadenoma} ($\geq$10 mm): may extend suprasellarly, compress optic chiasm, invade cavernous sinus
  \item \textbf{Dynamic contrast}: improves detection of microadenomas by showing differential enhancement between tumour and normal gland
  \item \textbf{Optic chiasm assessment}: essential when macroadenoma is present---compression requires urgent surgical referral
\end{itemize}

\subsection{Pituitary Apoplexy}
\begin{itemize}
  \item Sudden haemorrhage or infarction within a pituitary adenoma
  \item \textbf{MRI}: high T1 signal (blood), expanded sella, chiasmal compression
  \item \textbf{Clinical}: sudden headache, visual loss, ophthalmoplegia, adrenal crisis---neurosurgical emergency
\end{itemize}

\section{Pancreatic Neuroendocrine Tumours (pNET)}
\begin{itemize}
  \item \textbf{CT with pancreatic protocol} (arterial phase): hypervascular mass in pancreas
  \item \textbf{MRI}: superior contrast resolution; may detect small lesions missed on CT
  \item \textbf{Somatostatin receptor imaging}: Ga-68 DOTATATE PET/CT for well-differentiated pNETs (SSTR-positive); sensitivity $>$90\% for metastatic disease
  \item \textbf{FDG-PET/CT}: for high-grade poorly differentiated pNETs (SSTR-negative, FDG-avid)
  \item \textbf{Functional tumours}: insulinoma (small, often $<$2 cm), gastrinoma (often in duodenum), glucagonoma (larger at diagnosis)
\end{itemize}

\begin{pearlbox}[When to recommend MIBG scan]
I-123 MIBG (metaiodobenzylguanidine) is used for:
\begin{enumerate}
  \item Localising phaeochromocytoma (especially extra-adrenal, bilateral, or metastatic)
  \item Neuroblastoma staging
  \item Carcinoid tumour localisation (less sensitive than DOTATATE PET)
\end{enumerate}
MIBG is reserved for cases where anatomical imaging is inconclusive or for treatment planning (I-131 MIBG therapy).
\end{pearlbox}

\section{Interpretation}
Describe size, composition, enhancement, invasion, calcification, vascularity, lymph nodes, bone involvement, and functional tracer uptake. Imaging categories should be applied with the correct clinical and laboratory context.

\section{Safety}
Consider iodinated contrast, renal function, pregnancy, radiation, catecholamine-related risk, and procedural planning. Acute adrenal hemorrhage, pituitary apoplexy, airway compromise, and hormone-related emergencies require prompt communication.

\clearpage
